\documentclass[english, parskip=full]{scrartcl}
\usepackage[utf8]{inputenc}  % Kodierung der Datei
\usepackage[english]{babel}  % Sprache auf Deutsch 
\usepackage{color}
\usepackage{graphicx, gensymb}
\usepackage{amsfonts, cancel, amssymb}
\usepackage{amsmath, amsthm, array, esint}
\usepackage{booktabs, multirow}  % schönere Tabellen
\usepackage{caption}  % Anpassung der Captions
\usepackage{scrlayer-scrpage}    % Manipulation des Seitenstils
\pagestyle{scrheadings}  % Stil für die Seite setzen
\automark{section}  % Abschnittsnamen als Seitenbeschriftung 
\setheadsepline{.5pt}    % Kopzeile durch Linie abtrennen
\usepackage[hidelinks]{hyperref}  % Links und weitere PDF-Features
\usepackage{typearea, tikz, pgffor, verbatim}
\usetikzlibrary{calc, arrows.meta,arrows, graphs, quotes, positioning,angles}
\usepackage[cache=false, outputdir=build]{minted}
\usemintedstyle{vs}
\usepackage{placeins} % provides \FloatBarrier

\definecolor{bg}{rgb}{0.9,0.9,0.9}
\newcommand\numberthis{\addtocounter{equation}{1}\tag{\theequation}}
\usepackage{svg}
\usepackage{siunitx}
\usepackage{physics}
\usepackage{tensor}
\usepackage[compat=1.1.0]{tikz-feynman}
\renewcommand{\bra}{}
\newcommand{\kom}{}
\newcommand{\brak}{}
\renewcommand{\braket}{}
\newcommand{\intx}{}
\newcommand{\intr}{}
\newcommand{\kla}{}
\newcommand{\klam}{}
\newcommand{\klammer}{}
\newcommand{\oper}{}
\newcommand{\tecxt}{}
\newcommand{\Bra}{}
\newcommand{\Ket}{}
\renewcommand{\i}{\text{i}}
\newcommand{\e}{\text{e}}
\newcommand*{\Fourier}{\textsc{Fourier}\ }
\newcommand*{\F}{\mathcal{F}}
\newcommand*{\sinc}{\text{sinc\,}}
\newcommand{\conv}{\mathop{\scalebox{3.5}{\raisebox{-0.38ex}{\makebox[0.5\width][c]{$\ast$}}}}\limits}

\newcommand*{\titel}{On the Optimization Problem of Volume per Surface Area}
\newcommand*{\autor}{Joachim Schwardt}

\hypersetup{pdfauthor={\autor}, pdftitle={\titel}}

%\titlehead{titlehead}
\subject{Variational Calculus}
\title{\titel}
\author{\autor}
\date{\today}

\newcommand{\vol}{\text{Vol}}

\begin{document}

%\begin{titlepage}
\maketitle  % Titel setzen
%\tableofcontents  % Inhaltsverzeichnis setzen
%\end{titlepage}

\section{The Problem}
Consider an $N$-dimensional manifold $M$ with boundary $\partial M$.
We want to find the $M$ for which the ratio of $N$-dimensional volume $V := \vol_N(M)$ and $N-1$-dimensional ``area'' $A := \vol_{N-1}(\partial M)$ is maximised.
To that end we define
\begin{align}
\vol(M) &:= \intr\int_{\mathbb{R}^{N}} \Theta_M(x) \, \mathrm{d}^{N}x
\end{align}
where $\Theta_M(x) = 1$ for all $x\in M$ is the characteristic function.
Then the problem as stated above is equivalent to maximizing the ratio
\begin{align}
R &:= \frac{\vol_N(M)}{\vol_{N-1}(\partial M)}.
\end{align}
Without loss of generality we may now assume two things.
First, it is useful to fix the lower dimensional volume to unity, i.e. $A=1$, since the ratio is unaffected by that choice.
This also avoids any potential problems with $A=0$.
Second, we may assume $M$ to be convex. 
Without giving a formal proof, it should be clear how any concavities can be removed by both increasing $V$ while decreasing $A$.
This allows us to choose the ``center of mass'' of $M$ as our origin, while making sure that we can reach $\partial M$ with a straight line of unspecified length in any direction.
Then we may rewrite the ratio $R$ as
\begin{align}
R &= \intr\int_{\mathbb{R}^{N}} \Theta_M(x) \, \mathrm{d}^{N}x.
\end{align}
\subsection{General Spherical Coordinates}
We can define $N$-dimensional spherical coordinates by introducing angles $\phi\in [0,2\pi)$ and $\theta_1,\dots,\theta_{N-2} \in [0,\pi)$.
One can show that the volume element is given by
\begin{align}
\mathrm{d}^Nx &= r^{N-1} \sin(\theta_1)^{N-2} \sin(\theta_2)^{N-3} \dots \sin(\theta_{N-2}) \mathrm{d}r \mathrm{d}\phi \tecxt\text{d}\theta_1 \dots \tecxt\text{d}\theta_{N-2}.
\end{align}
Therefore, in spherical coordinates we can express $R$ using the function $r(\phi,\theta_1,\dots,\theta_{N-2})$ which is defined as the radial distance of $\partial M$ from the origin in an arbitrary direction specified by the set of angles.
We find
\begin{align*}
R &= \intx\int_{0}^{\pi} \dots \intx\int_{0}^{\pi} \intx\int_{0}^{2\pi}  \intx\int_{0}^{\infty} \Theta\kla\left( r(\phi,\theta_1,\dots,\theta_{N-2}) - r \right) \times\\
&\hspace{2cm} \times \, r^{N-1}\mathrm{d}r \mathrm{d}\phi \sin(\theta_1)^{N-2} \sin(\theta_2)^{N-3} \dots \sin(\theta_{N-2}) \tecxt\text{d}\theta_1 \dots \tecxt\text{d}\theta_{N-2} \numberthis = \\
&=  \int_{0}^{\pi} \intx\int_{0}^{2\pi} \dots \intx\int_{0}^{\pi} \frac{1}{N} r(\phi,\theta_1,\dots,\theta_{N-2})^N \, \mathrm{d}\phi \prod_{j=1}^{N-2} \sin(\theta_j)^{N-1-j} \mathrm{d}\theta_j. \numberthis
\end{align*}
\subsection{Lagrange Multipliers}
Since we fixed the area $A=1$ we need to include a corresponding term in our functional
\begin{align}
R &\rightarrow R + \lambda (A - 1).
\end{align}
Note that the $N-1$-dimensional volume element on the surface of an $N$-dimensional sphere of radius $r$ is identical to the angular part of $\tecxt\text{d}^Nx$.
Thus the area is given by
\begin{align}
A &= \int_{0}^{\pi} \intx\int_{0}^{\pi} \dots \intx\int_{0}^{2\pi} \sqrt{r^2 + \partial_\phi r(\dots)^2 + \partial_{\theta_1} r(\dots)^2 + \dots}^{N-1}  \mathrm{d}\phi \prod_{j=1}^{N-2} \sin(\theta_j)^{N-1-j} \mathrm{d}\theta_j.
\end{align}
The functional we need to maximize then becomes $R[r(\dots)] = \int \mathcal{R}(r(\dots), \partial r(\dots))$ where we defined
\begin{align}
\mathcal{R} &= \frac{1}{N} r(\phi,\theta_1 ,\dots,\theta_{N-2})^N + \lambda \sqrt{r^2 + \partial_\phi r(\dots)^2 + \partial_{\theta_1}r(\dots)^2 + \dots + \partial_{\theta_{N-2}} r(\dots)^2}^{N-1}.
\end{align}
FIXME: The surface area formula feels wrong, and I also do not know how to derive it -- this was more of a guess...
\section{Solution to the Optimization Problem}
\subsection{Case $N=2$}
In two dimensions the functional is $R[r(\phi)] = \intx\int_{0}^{2\pi} \mathcal{R}(r(\phi), r'(\phi)) \, \mathrm{d}\phi$ with
\begin{align}
\mathcal{R} &= \frac{1}{2} r(\phi)^2 + \lambda\sqrt{r^2 + r'(\phi)^2}.
\end{align}
Since there is no explicit dependence on $\phi$ the Euler-Lagrange equations become
\begin{align}
\mathcal{R} - r'(\phi) \frac{\partial \mathcal{R}}{\partial r'(\phi)} &= \tecxt\text{const} =: c.
\end{align}
In our case we find
\begin{align}
c &= \frac{r^2}{2} + \lambda\sqrt{r^2 + r'^2} - \lambda \frac{r'^2}{\sqrt{r^2 + r'^2}}.
\end{align}
Multiplying by the square root yields
\begin{align}
\kla\left( c - \frac{r^2}{2} \right) \sqrt{ r^2 + r'^2 } &= \lambda r^2.
\end{align}
Squaring and rearranging gives
\begin{align}
r'^2 &= \frac{\lambda^2 r^4}{\kla\left( c - \frac{r^2}{2} \right)^2} - r^2 =: f(r).
\end{align}
The solution to this differential equation is given by 
\begin{align}
\phi + \tecxt\text{const} &= \intx\int_{ }^{ } \frac{1}{\sqrt{f(r)}} \, \mathrm{d}r = \intx\int_{ }^{ } \frac{1}{r} \kla\left( \frac{\lambda^2r^2}{c^2 - cr^2 + \frac{r^4}{4}} - 1 \right)^{-\frac{1}{2}} \, \mathrm{d}r = \\
&= \frac{1}{2} \intx\int_{ }^{ } \frac{1}{u} \kla\left( \frac{\lambda^2}{c^2 - cu + \frac{u^2}{4}} - 1 \right)^{-\frac{1}{2}} \, \mathrm{d}u = \frac{1}{2} \intx\int_{ }^{ } \frac{1}{u} \sqrt{\frac{c^2 - cu + \frac{u^2}{4}}{\lambda^2 - c^2 + cu - \frac{u^2}{4}}} \, \mathrm{d}u
\end{align}



%\begin{thebibliography}{9}
%\bibitem{example} description, \url{https://www.exmapleurl}, day.\,month~2021.
%\end{thebibliography}

\end{document}